% bianchi_extension_opus.tex
% Companion / appendix to algebraic_arrow.tex --- attempt to extend
% Algebraic Weyl Curvature Hypothesis (T1+T2+T3) from FRW to Bianchi.
%
% Author: Kevin Remondiere with max-effort theorem attempt by Opus 4.7.
% Date: 2026-05-02
% Sympy verification: /tmp/bianchi_extension_opus.py

\documentclass[11pt]{article}
\usepackage[utf8]{inputenc}
\usepackage[T1]{fontenc}
\usepackage[margin=1in]{geometry}
\usepackage{amsmath,amssymb,amsthm,mathtools,hyperref,xcolor,booktabs}
\hypersetup{colorlinks=true,linkcolor=blue,citecolor=blue,urlcolor=blue}

\newtheorem{theorem}{Theorem}
\newtheorem{proposition}[theorem]{Proposition}
\newtheorem{lemma}[theorem]{Lemma}
\newtheorem{corollary}[theorem]{Corollary}
\theoremstyle{definition}
\newtheorem{remark}[theorem]{Remark}

\newcommand{\R}{\mathbb{R}}
\newcommand{\C}{\mathbb{C}}
\newcommand{\HH}{\mathcal{H}}
\renewcommand{\AA}{\mathcal{A}}
\newcommand{\OO}{\mathcal{O}}
\newcommand{\Ad}{\operatorname{Ad}}
\newcommand{\BI}{\mathrm{BI}}
\newcommand{\BV}{\mathrm{BV}}
\newcommand{\BIX}{\mathrm{BIX}}
\newcommand{\FRW}{\mathrm{FRW}}
\newcommand{\Mink}{\mathrm{M}}
\newcommand{\BKL}{\mathrm{BKL}}
\newcommand{\HL}{\mathrm{HL}}
\newcommand{\BFV}{\mathrm{BFV}}

\title{Algebraic Weyl Curvature Hypothesis: Bianchi extension attempt\\
  \large Perturbative success, non-perturbative obstructions}
\author{Kevin Remondi\`ere\\
  \small with max-effort theorem attempt by Opus 4.7 (1M ctx)}
\date{2026-05-02 (companion to \texttt{algebraic\_arrow.tex})}

\begin{document}
\maketitle

\begin{abstract}
We attempt to extend the Algebraic Weyl Curvature Hypothesis (T1+T2+T3 of
\cite{algebraic_arrow}, proved modulo the Hadamard-folium gap for FRW) to
Bianchi cosmologies. Bianchi types I, V, IX are not conformally flat:
sympy verification shows the vacuum Kasner Weyl tensor is non-zero and
divergent at the singularity, e.g.\
$C_{txtx} = -\tfrac{4}{9 t^{8/3}}$ for the standard exponents
$(p_1,p_2,p_3) = (-1/3,2/3,2/3)$. Penrose's WCH is therefore non-trivial
on Bianchi backgrounds, in contrast to FRW (where $C_{abcd}\equiv 0$ by
symmetry). Five strategies are evaluated:

(S1) Direct GNS: PARTIAL --- Hadamard state exists (Pereira--Bezerra de Mello--Brevik
\cite{Pereira2018}, Them--Brum \cite{ThemBrum2013} states-of-low-energy on
inhomogeneous spacetimes), but no analog of Hislop--Longo modular flow because
no conformal Killing field generates a M\"obius flow on causal diamonds.

(S2) Mini-superspace reduction: FAILS --- Connes--Takesaki crossed product
with the Bianchi Lie group $G$ collapses the type structure to type I$_\infty$.

(S3) BKL bridge \`a la Hartnoll--Yang \cite{HartnollYang2025,
ClerckHartnoll2025}: STRUCTURAL ANALOGY ONLY --- the Hecke
algebra of automorphic forms is type I$_\infty$ abelian,
not type III.

(S4) Tod isotropic singularity boundary: PARTIAL --- recovers FRW T2 exactly
on the Weyl-vanishing class, gives no extension to Weyl-divergent class.

(S5) Perturbative anisotropy lift around FRW: \textbf{SUCCEEDS at first
order} via the Connes 1973 cocycle invariance theorem (\textbf{Proposition B.1}
below). This is the BEST RESULT obtained.

The unconditional non-perturbative extension hits three open problems
(\S\ref{sec:obstructions}); estimated effort to crack any one of them is
6--12 months; full Bianchi I/V/IX extension is a 2--5 year program with
serious risk of fundamental obstruction.
\end{abstract}

\section{Bianchi backgrounds: Weyl tensor non-vanishing}\label{sec:bianchi}

We work in synchronous time $t$ throughout. Sympy verification (script
\texttt{/tmp/bianchi\_extension\_opus.py}) confirms the following:

\begin{table}[h]
\centering
\begin{tabular}{lll}
\toprule
Type & Metric (synchronous) & Weyl tensor at singularity \\
\midrule
\textbf{I} (Kasner) &
  $-\mathrm dt^2 + \sum_i t^{2 p_i}\,\mathrm dx_i^2$ &
  $C_{txtx} = -\tfrac{4}{9 t^{8/3}}$, $\to\infty$\\
& vacuum: $\sum p_i = \sum p_i^2 = 1$ &
  (sympy C1 verified) \\
\textbf{V} & $-\mathrm dt^2 + t^{2p_1}\mathrm dx^2$ & $C_{txtx}=\tfrac{31}{864\,t}$\\
& $+\,t^{2p_2} e^{2x}\mathrm dy^2 + t^{2p_3}e^{2x}\mathrm dz^2$ & (sympy C3 verified)\\
\textbf{IX} (Mixmaster) &
  $-\mathrm dt^2 + \sum_i \ell_i(t)^2 (\sigma^i)^2$ &
  non-zero, BKL chaotic\\
& $\sigma^i$ Maurer-Cartan on SU(2) &
  bouncing near $t\!\to\!0$ \\
\bottomrule
\end{tabular}
\end{table}

\noindent\textbf{Critical obstruction (sympy C2).}
Weyl is a conformal invariant in $d=4$:
$C^a{}_{bcd}[\Omega^2 g] = C^a{}_{bcd}[g]$ pointwise. Since
$C_{abcd}[\text{Mink}]\equiv 0$ but $C_{abcd}[\text{Kasner}]\not\equiv 0$,
\emph{no global conformal factor $\Omega(x)$} maps Bianchi I to Minkowski.
The conformal-pullback unitary $U$ of \cite[Prop.~3.3]{frw_note} has no
Bianchi I analog. This is the fundamental reason all five strategies
struggle.

\section{Strategy-by-strategy verdicts}\label{sec:strategies}

\subsection{Strategy 1: Direct GNS construction (Bianchi I)}

\textbf{Verdict: PARTIAL.} Pereira--Bezerra de Mello--Brevik
\cite{Pereira2018} construct a preferred adiabatic Hadamard state
$\Omega_\BI$ on Bianchi I/Kasner via mode decomposition along the
abelian $\R^3$ symmetry; Them--Brum \cite{ThemBrum2013} extend the
states-of-low-energy framework from RW to inhomogeneous spacetimes
(includes Bianchi I).

\smallskip\noindent\emph{What works.}
The local algebra $\AA(D)_\BI$ for a doubly-bounded comoving causal
diamond $D$ admits a cyclic-separating vector $\Omega_\BI$ in its GNS
representation. Microlocal spectrum condition
\cite{BrunettiFredenhagenVerch2003} holds. Hence $\AA(D)_\BI$ is
plausibly a type III$_1$ factor (by general arguments of
Buchholz--Verch on Hadamard states satisfying the split property).

\smallskip\noindent\emph{What fails.}
There is no analog of the Hislop--Longo modular flow because Bianchi I
has \emph{no conformal Killing field generating a M\"obius flow on
causal diamonds}. The only Killing vectors are the 3 spatial
translations $\partial_{x^i}$, which act tangentially to Cauchy slices
and cannot generate the modular automorphism (the modular generator
$K_\BI$ is essentially boost-like on the Cauchy slice, not translational).
Without an explicit modular-flow identification, the Connes spectrum
$\Gamma(\sigma^\BI)$ cannot be computed by the technique of
\cite{frw_note}.

\subsection{Strategy 2: Symmetric reduction (mini-superspace)}

\textbf{Verdict: FAILS.}
Reduce QFT on Bianchi I via Connes--Takesaki crossed product with the
Bianchi Lie group $G\subset\mathrm{Diff}(M)$:
$\AA(D)_\BI \rtimes G$. But $G$ is a finite-dim Lie group (3-dim
abelian for Bianchi I), so the reduction collapses the QFT to
mini-superspace QM on $G\backslash\Sigma\times\R$. This is
\emph{type I$_\infty$} (separable Hilbert space, abelian center),
losing the type III/II structure entirely. This is a known obstruction
in mini-superspace quantization: integrating out
infinitely many modes destroys the von Neumann type structure of the
underlying QFT.

\subsection{Strategy 3: BKL bridge \`a la Hartnoll--Yang}

\textbf{Verdict: STRUCTURAL ANALOGY ONLY.}
Hartnoll--Yang \cite{HartnollYang2025} and Clerck--Hartnoll
\cite{ClerckHartnoll2025} establish that BKL Bianchi IX dynamics map
onto a particle in a fundamental domain of $PSL(2,\Z)$ (resp.\
$PSL(2,\OO)$ for 5d, $\OO$ a quadratic imaginary order), with
Wheeler--DeWitt wavefunctions equal to automorphic L-functions on the
critical axis. Modular invariance of states is enforced.

\smallskip\noindent\emph{What works.}
The natural Hilbert space $\HH_\BKL = L^2(F, \mathrm{dvol}_\mathrm{hyp})$
with $F$ the fundamental domain of $PSL(2,\Z)$ admits Maass cusp forms
as cyclic vectors for the Hecke operator algebra
$\langle T_p\rangle_p$.

\smallskip\noindent\emph{What fails.}
The Hecke algebra is \emph{commutative} ($T_p T_q = T_q T_p$ for all
distinct primes), so its von Neumann closure is type I$_\infty$
abelian $L^\infty(\mathrm{spec})$, not type III or II. The
cyclic-separating-vector + Connes-spectrum framework requires a
non-abelian type III/II factor. Bridge gives a suggestive parallel
(modular-invariant ground state $\leftrightarrow$ Bunch--Davies cyclic
vector) but no direct theorem.

\subsection{Strategy 4: Tod isotropic singularity boundary condition}

\textbf{Verdict: PARTIAL --- gives a CLASSIFICATION not an extension.}
Tod's isotropic singularity theorem \cite{Tod1991, GoodeWainwright1985}
characterizes singularities at which the Weyl tensor is dominated by the
Ricci tensor in a Penrose-conformal sense, so that after rescaling by a
Penrose factor $\widehat\Omega$ with $\widehat\Omega \to 0$ at the
singularity, the rescaled metric is regular and $C_{abcd}\to 0$ at the
conformal boundary.

In our framework: \emph{isotropic singularities reduce to the FRW T2
theorem} (the conformal pullback works exactly because Weyl-vanishing
is the defining condition). This is a consistency check (and confirms
that our framework is in the same spirit as Penrose's WCH:
``Weyl-vanishing at sing.'' is exactly when our T2 obstruction is
removable). It provides no extension to the Weyl-divergent class
(Kasner, BKL, generic Bianchi).

\subsection{Strategy 5: Perturbative anisotropy lift around FRW}

\textbf{Verdict: SUCCEEDS at first order.} See \S\ref{sec:propB1}.

\section{Proposition B.1 (perturbative Bianchi I lift)}\label{sec:propB1}

\begin{proposition}[Perturbative type II$_\infty$ for near-FRW Bianchi I]\label{prop:B1}
Let $\bigl(M, g_\BI(\varepsilon)\bigr)$ be a one-parameter family of
Bianchi I spacetimes with $g_\BI(0) = g_\FRW$ (radiation- or
matter-dominated, spatially flat $k=0$) and small anisotropy parameter
$\varepsilon\in[0,\varepsilon_0)$:
\begin{equation}\label{eq:perturb-metric}
\mathrm{d}s^2_{\BI(\varepsilon)} \;=\; -\mathrm dt^2
   \,+\, \sum_{i=1}^3 a(t)^2 \bigl(1 + \varepsilon\,\delta_i(t)\bigr)^2\,\mathrm dx^i\,\mathrm dx^i,
\qquad \sum_i \delta_i \equiv 0
\end{equation}
(traceless anisotropy, volume preserving to first order). Let $\varphi$
be the conformally coupled massless scalar in $d=4$, and
$\Omega_{\BI(\varepsilon)}$ the adiabatic Hadamard / states-of-low-energy
ground state of \cite{Pereira2018, ThemBrum2013}. Let $D\subset M$ be
a doubly-bounded comoving causal diamond with $0<t_i<t_f<\infty$. Then
there exists $\varepsilon_0 = \varepsilon_0(D)>0$ such that for all
$\varepsilon\in[0,\varepsilon_0)$:
\begin{enumerate}
\item $\AA(D)_{\BI(\varepsilon)}$ is a type III$_1$ factor in the GNS rep
  of $\Omega_{\BI(\varepsilon)}$.
\item $\Gamma\bigl(\sigma^{\BI(\varepsilon)}\bigr) = \R$.
\item $\AA(D)_{\BI(\varepsilon)} \rtimes_\sigma \R$ is a type II$_\infty$
  factor.
\end{enumerate}
\end{proposition}

\begin{proof}
For $\varepsilon=0$ this is Theorem~3.5 of \cite{frw_note} via the
conformal-pullback unitary $U$. For $\varepsilon\in(0,\varepsilon_0)$:

\textbf{Step 1 (BFV local quasi-equivalence).} By Brunetti--Fredenhagen--Verch
\cite{BrunettiFredenhagenVerch2003} Thm.~4.1 + Hollands--Wald
\cite{HollandsWald2001} Sec.~5, the GNS reps $\pi_0:=\pi_{\Omega_\FRW}$ and
$\pi_1:=\pi_{\Omega_{\BI(\varepsilon)}}$ on the local algebra $\AA(D)$ are
quasi-equivalent: there exists a $W^*$-isomorphism $\pi_1(\AA(D))''
\xrightarrow{\sim}\pi_0(\AA(D))''$ on the same closure $\bar\AA(D)$.
The threshold $\varepsilon_0(D)$ is the BFV quasi-equivalence radius:
estimate $\varepsilon_0\sim \mathrm{(curvature\ radius\ of\ }D)/
\mathrm{(diameter\ of\ }D)$, well within CMB anisotropy bounds
($\sim 10^{-5}$ on largest scales) for cosmologically realistic $D$.

\textbf{Step 2 (Connes cocycle).} By Connes \cite{Connes1973}
Th\'eor\`eme~1.2.4 (Radon--Nikodym for modular automorphisms): for
faithful normal states $\omega_0,\omega_1$ on the same vN algebra $M$,
the modular flows differ by a $\sigma$-cocycle $u_t :=
(D\omega_1\!:\!D\omega_0)_t$ inside $M$, hence
\begin{equation}\label{eq:connes-cocycle}
\sigma_t^{\omega_1} \;=\; \Ad u_t \circ \sigma_t^{\omega_0}.
\end{equation}
Apply with $M = \pi_0(\AA(D))''$, $\omega_0 = $ Bunch--Davies state on
FRW, $\omega_1 = $ adiabatic Hadamard state on $\BI(\varepsilon)$ pushed
through the BFV isomorphism of Step 1.

\textbf{Step 3 (Connes spectrum invariance).} By Connes
\cite{Connes1973} Cor.~1.2.4, the Connes spectrum is invariant under
inner perturbation of the modular automorphism by a $\sigma$-cocycle
inside $M$:
\begin{equation*}
\Gamma(\sigma^{\omega_1}) \;=\; \Gamma(\sigma^{\omega_0}) \;=\;
\Gamma(\sigma^\HL) \;=\; \R
\end{equation*}
(last equality by Hislop--Longo \cite{HislopLongo1982} via FRW Theorem
3.5).

\textbf{Step 4 (type classification).} By Connes--Takesaki, type III$_1$
crossed by its modular automorphism group with full Connes spectrum
$\Gamma=\R$ is type II$_\infty$.

The type III$_1$ factoriality of $\AA(D)_{\BI(\varepsilon)}$ (item 1) is
preserved by the $W^*$-isomorphism of Step 1 from the FRW type III$_1$
factor $\AA(D)_\FRW$ (which is type III$_1$ by our \cite[Thm.~3.5]{frw_note}
via Hislop--Longo).
\end{proof}

\noindent\textbf{Sympy verification (script
\texttt{/tmp/bianchi\_extension\_opus.py} check C4).} For radiation FRW
$a(t)=\sqrt t$ and the perturbed metric \eqref{eq:perturb-metric}, the
Riemann component $R^t{}_{xtx}$ is computed as a Taylor series in
$\varepsilon$:
\begin{equation*}
R^t{}_{xtx}\bigl(\varepsilon\bigr)
   \;=\; -\frac{1}{4t} \;-\; \frac{\delta_1\,\varepsilon}{2t}\;+\;O(\varepsilon^2),
\end{equation*}
with the $\varepsilon^0$ term equal to the FRW Riemann (Weyl-vanishing
part) and the $\varepsilon^1$ term carrying the leading anisotropic
Weyl correction. This confirms that the perturbative expansion is
well-defined and that $g_\BI(\varepsilon)$ is a smooth one-parameter
family with $g_\BI(0)=g_\FRW$.

\section{Critical obstructions to non-perturbative extension}\label{sec:obstructions}

The unconditional Bianchi extension requires resolving \emph{at least
one} of three open problems:

\paragraph{Problem 1: Non-conformal modular-flow identification.}
For Bianchi I (and a fortiori V/IX), the modular flow of $\Omega_\BI$
on a comoving diamond is NOT generated by any conformal Killing field.
Open: characterize $\sigma^{\BI}_t$ explicitly, or sidestep with modular
nuclearity bounds (Buchholz--D'Antoni--Longo) or Buchholz--Mund--Summers
transplantation \cite{BuchholzMundSummers2000} of the wedge-net CGMA
+ modular stability without explicit modular flow.

\paragraph{Problem 2: Hadamard uniqueness for anisotropic singularities.}
Hollands--Wald \cite{HollandsWald2001} + BFV
\cite{BrunettiFredenhagenVerch2003} give Hadamard uniqueness on globally
hyperbolic spacetimes; the singularity is a true boundary, not in the
manifold. For BKL singularities, the relevant uniqueness theorem may
not exist --- Damour--Henneaux--Nicolai cosmological billiards
\cite{DHN2002} suggest the asymptotic behavior is Kac--Moody-controlled
($E_{10}$ for 11d supergravity; $A_1^{++}$ for vacuum 4d Bianchi IX),
which is far richer than the Hadamard parametrix structure.

\paragraph{Problem 3: Type-classification of $\AA(D)_\BI$.}
This is the Bianchi analog of Hislop--Longo \cite{HislopLongo1982}.
Status: genuinely open; no published result computes the von Neumann
type for the conformally coupled scalar in vacuum Kasner.

\section{Recommendation and effort estimate}\label{sec:recommendation}

\begin{itemize}
\item \textbf{For \texttt{algebraic\_arrow.tex}}: ADD a Section 6
  (``Perturbative anisotropy lift'') stating Proposition~\ref{prop:B1}
  with the cocycle proof above (~2 pages incl.\ sympy verification).
  This converts the ``Bianchi extension is open'' caveat in Sec.~5 of
  that note from a pure caveat to a positive first-order result.

\item \textbf{For a standalone paper}: NOT recommended yet. Better
  option: a methodological note ``Why FRW conformal-pullback is
  fundamentally conformally flat: Bianchi I obstructions and the
  perturbative cocycle workaround'' for \emph{Class.\ Quant.\ Grav.}
  Comments and Replies (~5 pages, ~3-4 weeks effort).

\item \textbf{Effort estimate} for the unconditional extension:
  \begin{itemize}
  \item Crack one of Problems 1--3: 6--12 months for an expert in
    algebraic QFT in curved spacetime.
  \item Full Bianchi I/V/IX unconditional extension: 2--5 years, with
    serious risk of fundamental obstruction (especially BKL Bianchi IX
    where chaos breaks the BFV folium structure).
  \end{itemize}
\end{itemize}

\bibliographystyle{plain}
\begin{thebibliography}{99}
\bibitem{algebraic_arrow} K.~Remondi\`ere, with Opus 4.7,
\emph{Algebraic Weyl Curvature Hypothesis: proof attempt for the FRW
Type II$_\infty$ framework},
\texttt{notes/path\_realiste\_2026\_05\_03/algebraic\_arrow.tex} (2026-05-02).

\bibitem{frw_note} K.~Remondi\`ere, \emph{Type II$_\infty$ for the
gravitational algebra of a comoving observer in radiation-dominated FRW
via conformal pullback},
\texttt{paper/frw\_typeII\_note/frw\_note.tex} (2026-05-02).

\bibitem{HartnollYang2025} S.~A.~Hartnoll, M.~Yang,
\emph{The Conformal Primon Gas at the End of Time},
\texttt{arXiv:2502.02661} (2025).

\bibitem{ClerckHartnoll2025} L.~De Clerck, S.~A.~Hartnoll,
\emph{Wheeler--DeWitt wavefunctions for 5d BKL dynamics, automorphic
L-functions and complex primon gases}, \texttt{arXiv:2507.08788} (2025).

\bibitem{DHN2002} T.~Damour, M.~Henneaux, H.~Nicolai,
\emph{Cosmological Billiards}, Class.\ Quant.\ Grav.\ \textbf{20}
(2003) R145, \texttt{arXiv:hep-th/0212256}.

\bibitem{Tod1991} K.~P.~Tod, \emph{Isotropic cosmological
singularities: other matter models}, Class.\ Quant.\ Grav., 1990s
series.

\bibitem{GoodeWainwright1985} S.~W.~Goode, J.~Wainwright,
\emph{Isotropic singularities in cosmological models}, Class.\ Quant.\
Grav.\ \textbf{2} (1985) 99.

\bibitem{Pereira2018} G.~Pereira, E.~R.~Bezerra de Mello, I.~Brevik,
\emph{Quantum fields in Bianchi type I spacetimes: The Kasner metric},
Phys.\ Rev.\ D \textbf{98} (2018) 104054, \texttt{arXiv:1811.00993}.

\bibitem{ThemBrum2013} K.~Them, M.~Brum, \emph{States of Low Energy in
Homogeneous and Inhomogeneous, Expanding Spacetimes},
\texttt{arXiv:1302.3174} (2013).

\bibitem{BrunettiFredenhagenVerch2003} R.~Brunetti, K.~Fredenhagen,
R.~Verch, \emph{The generally covariant locality principle}, Comm.\
Math.\ Phys.\ \textbf{237} (2003) 31, \texttt{arXiv:math-ph/0112041}.

\bibitem{HollandsWald2001} S.~Hollands, R.~M.~Wald, \emph{Local Wick
polynomials and time ordered products of quantum fields in curved
spacetime}, Comm.\ Math.\ Phys.\ \textbf{223} (2001) 289,
\texttt{arXiv:gr-qc/0103074}.

\bibitem{Connes1973} A.~Connes, \emph{Une classification des facteurs
de type III}, Ann.\ Sci.\ ENS \textbf{6} (1973) 133.

\bibitem{HislopLongo1982} P.~D.~Hislop, R.~Longo, \emph{Modular
structure of the local algebras associated with the free massless
scalar field theory}, Comm.\ Math.\ Phys.\ \textbf{84} (1982) 71.

\bibitem{BuchholzMundSummers2000} D.~Buchholz, J.~Mund, S.~J.~Summers,
\emph{Transplantation of local nets and geometric modular action on
Robertson--Walker space-times}, \texttt{arXiv:hep-th/0011237}.

\bibitem{Penrose2010} R.~Penrose, \emph{Cycles of Time: An
Extraordinary New View of the Universe}, Bodley Head 2010; see also
\texttt{arXiv:1011.3706}.
\end{thebibliography}

\end{document}
